% SGA 3, Expose VI_A, Loop 1 English reconstruction.
% Scope: paragraph 2.6 through the end of Proposition 2.6.4;
% Theorem 2.6.5 is excluded.
% Authority: Polo--Gille Exp6A-13oct24.pdf, component-local pages 15--18
% (combined-reader pages 318--321; printed pages 307--310).
% OCR and the Reinhold English Markdown are comparison witnesses only.

\medskip
\phantomsection\label{sec:sga3-VIA-2.6}%
\noindent\SGAInlineRefTarget{sga3:vi:paragraph:2-6}\textbf{2.6. 补充.---}
本段为新增内容, 取自 [Per75], II, \S\S~1--2,
并包含 O.~Gabber 所作的补充.%
\footnote{\label{note:sga3-VIA-31}编者注: 这些结果由 O.~Gabber 告知我们,
特别是\SGARefLink{prop:sga3-VIA-2.6.6}{2.6.6}, 它在第 VI\textsubscript{B} 篇
\hyperref[sec:sga3-VIB-5]{第~5~节}中起重要作用.}
本段表明, 前面的结果对域 \(k\) 上的每个群概形 \(G\) 都成立.
(这将用于第 VI\textsubscript{B} 篇的\hyperref[sec:sga3-VIB-5]{第~5~节}、第 6 节及第 7 节.)

固定域 \(k\). 先给出下述引理 (同上, II, 2.1.1),
它未在 EGA~IV\textsubscript{2}, \S~4.4 中明确出现
(不过, 或许可从同上 \S~4.4.1 开头的字里行间读出此结论).

\medskip
\phantomsection\label{lem:sga3-VIA-2.6.0}%
\noindent\textbf{引理 2.6.0.}
\textit{设 \(X\) 为不可约 \(k\)-概形, \(K\) 为 \(k\) 的扩张,
\(X'\) 为 \(X_K\) 的一个不可约分支. 投影 \(X'\to X\) 为满射.}

\noindent\textit{证明.}
设 \(B\) 为 \(K\) 在 \(k\) 上的超越基, 并令 \(L=k(B)\subset K\).
由 EGA~IV\textsubscript{2}, 4.3.2 和 4.4.1,
\(X_L\) 不可约, 且 \(X'\) 支配 \(X_L\).
故 \(X'\to X_L\) 是整且支配的, 从而满射.
由于 \(X_L\to X\) 满射 (同上, 4.4.1), \(X'\to X\) 也满射.

在\hyperref[sga3:vi:paragraph:2-6]{2.6} 的其余部分, 固定 \(k\)-群概形 \(G\),
以及 \(G\) 在 \(k\)-概形 \(X\) 上的作用 \(\mu:G\times X\to X\),
使之满足下述条件:
\[
  \tag{\(\star\)}\label{sga3:vi:lemma:2-6-0:equation:star}
  \Phi:G\times X\longrightarrow X\times X,\qquad
  (g,x)\longmapsto(gx,x)
  \quad\text{为满射.}
\]
(特别地, 当 \(G\) 通过左平移作用于自身时, 此条件成立.)
为简便起见, 我们将说: ``设 \(X\) 为满足 \((\star)\) 的 \(G\)-概形.''
最后, 用 \(k'\) 表示 \(k\) 的完美闭包.

\medskip
\phantomsection\label{prop:sga3-VIA-2.6.1}%
\noindent\textbf{命题 2.6.1.}
\textit{设 \(X\) 为满足 \((\star)\) 的 \(G\)-概形.}
\begin{enumerate}
\item[\SGAInlineRefTarget{prop:sga3-VIA-2.6.1:item:i}(i)] \textit{\(X\) 在 \(k\) 上几何逐点不可约: 对 \(k\) 的每个扩张 \(K\),
每个点 \(x\in X_K\) 属于 \(X_K\) 的唯一一个不可约分支.}
\item[\SGAInlineRefTarget{prop:sga3-VIA-2.6.1:item:ii}(ii)] \textit{\((X_{k'})_{\mathrm{red}}\) 的每个局部环都正规.}
\item[\SGAInlineRefTarget{prop:sga3-VIA-2.6.1:item:iii}(iii)] \textit{设 \(\eta\) 为 \((X_{k'})_{\mathrm{red}}\) 的一个极大点,
令 \(C=\overline{\{\eta\}}\), 并设 \(L\) 为 \(k\) 在 \(\kappa(\eta)\) 中的代数闭包.
则 \(C\) 是 \(L\)-概形, 且在 \(L\) 上几何不可约;
即对 \(L\) 的每个扩张 \(K\), \(C\otimes_LK\) 不可约.}
\item[(iv)] \textit{特别地, 若 \(x\) 为 \(X\) 的有理点,
则 \(X\) 中包含 \(x\) 的不可约分支在 \(k\) 上几何不可约.}
\end{enumerate}

\noindent\textit{证明.}
\SGARefLink{prop:sga3-VIA-2.6.1:item:i}{(i)} 由于假设 \((\star)\) 在每个基变换 \(k\to K\) 下保持,
只需证明每个 \(x\in X\) 属于 \(X\) 的唯一一个不可约分支.
由于 \(\operatorname{Spec}(k')\to\operatorname{Spec}(k)\) 是普遍同胚,
还可假设 \(k\) 完美. 进而可假设 \(G\) 和 \(X\) 约化.

设 \(\eta\) 为 \(X\) 的极大点, \(z\) 为 \(Z=\overline{\{\eta\}}\) 的任意一点.
由于 \(X\) 约化, \(\mathcal O_{X,\eta}=\kappa(\eta)\);
又由于 \(k\) 完美, 对 \(k\) 的每个扩张 \(K\),
\(\kappa(\eta)\otimes_kK\) 正规 (参见 EGA~IV\textsubscript{2}, 6.14.2).
因此 \(X_K\) 中位于 \(\eta\) 上方的每个点都正规.

由于 \(\Phi\) 满射, 存在 \(G\times X\) 的一点 \(\gamma\),
使 \(\Phi(\gamma)\) 的两个投影为 \(z\) 和 \(\eta\).
令 \(K=\kappa(\gamma)\). 于是存在有理点 \(g\in G_K\) 和 \(\eta'\in X_K\),
其中 \(\eta'\) 位于 \(\eta\) 上方, 使得 \(z'=g\eta'\) 位于 \(z\) 上方.
由前一段, \(\eta'\) 为 \(X_K\) 的正规点, 从而 \(z'\) 也是正规点.
由于 \(\pi:X_K\to X\) 平坦, 可知 \(z=\pi(z')\) 为 \(X\) 的正规点
(参见 EGA~IV\textsubscript{2}, 2.1.13). 这就证明了 (ii) 和 (i).

现在, \SGARefLink{prop:sga3-VIA-2.6.1:item:iii}{(iii)} 的第一个断言由
\SGARefLink{prop:sga3-VIA-2.6.1:item:ii}{(ii)} 得到.
由于 \(L\) 在 \(\kappa(\eta)\) 中代数闭,
由 EGA~IV\textsubscript{2}, 4.5.9 可知 \(C\) 在 \(L\) 上几何不可约.

最后, 若 \(X\) 有有理点 \(x\), 则由\SGARefLink{prop:sga3-VIA-2.6.1:item:iii}{(iii)}
可知 \(L=k\), 故包含 \(x\) 的不可约分支 \(C\) 在 \(k\) 上几何不可约.\begingroup
\renewcommand{\thefootnote}{\fnsymbol{footnote}}%
\footnote[1]{英译者注: 法文底本在此结论中印作 \(X\), 随后又将点 \(x\)
换为未定义的符号 \(e\) 和 \(e_K\).
陈述及论证分别需要分支 \(C\) 以及点 \(x\) 和 \(x_K\).}%
\endgroup
也可如下直接看出此结论 (参见 [Per75], II, \SGARefLink{par:sga3-VIA-2.1}{2.1}).
设 \(C\) 为 \(X\) 中包含 \(x\) 的不可约分支, \(K\) 为 \(k\) 的扩张.
概形 \(X_K\) 在 \(x\) 上方有唯一一点 \(x_K\),
由\SGARefLink{prop:sga3-VIA-2.6.1:item:i}{(i)}, \(x_K\) 属于 \(X_K\) 的唯一一个不可约分支 \(C'\).
另一方面, 由\hyperref[lem:sga3-VIA-2.6.0]{2.6.0},
\(C_K\) 的每个不可约分支都包含 \(x_K\), 因而等于 \(C'\).

\medskip
\phantomsection\label{not:sga3-VIA-C0}%
\noindent\textbf{记号.---}
暂用 \(C^0\) 表示 \(G\) 的约化闭子概形,
其底空间为 \(G\) 中包含单位元 \(e\) 的唯一不可约分支.

\medskip
\phantomsection\label{cor:sga3-VIA-2.6.2}%
\noindent\textbf{推论 2.6.2.}
\textit{\(C^0\) 在 \(k\) 上几何不可约, 且在集合意义下对群运算稳定;
即 \((C^0_{k'})_{\mathrm{red}}\) 为 \((G_{k'})_{\mathrm{red}}\) 的子群.
因此 \(C^0\) 拟紧.}

\noindent\textit{证明.}
由\hyperref[prop:sga3-VIA-2.6.1]{2.6.1}, \(C^0\) 在 \(k\) 上几何不可约,
故 \(C^0\times C^0\) 不可约, 且
\[
  \nu(C^0\times C^0)\subset C^0,
\]
其中 \(\nu\) 表示态射 \((g,h)\mapsto gh^{-1}\).
由于 \(\operatorname{Spec}k'\to\operatorname{Spec}k\) 为普遍同胚,
对 \(C^0_{k'}\) 也有同样结论, 进而对
\[
  H=(C^0_{k'})_{\mathrm{red}}.
\]
也有同样结论. 由于 \(H\times_{k'}H\) 约化, \(\nu\) 诱导态射
\[
  H\times_{k'}H\longrightarrow H;
\]
故 \(H\) 为 \((G_{k'})_{\mathrm{red}}\) 的子群.
因此, 由\hyperref[sga3:VIA:0.5.1]{0.5.1}, \(H\) 拟紧, 从而 \(C^0\) 拟紧.

\medskip
\phantomsection\label{rem:sga3-VIA-2.6.3}%
\noindent\textbf{回顾 2.6.3.---}
设 \(Y\) 为概形. 回顾 (EGA~0\textsubscript{III}, 9.1.1),
\(Y\) 的子集 \(E\) 称为\emph{相对拟紧的 (retrocompact)},
若包含映射 \(E\hookrightarrow Y\) 拟紧.
此外, 由 EGA~IV\textsubscript{1}, 1.9.5(v) 和 1.10.1,
若 \(U\) 为 \(Y\) 的相对拟紧开子集, 则其闭包 \(\overline U\)
为点 \(y\in U\) 的闭包 \(\overline{\{y\}}\) 之并.
当然, 只需让 \(y\) 遍历 \(U\) 的极大点, 即 \(Y\) 中包含于 \(U\) 的极大点.

因此, 若 \(Y\) 拟分离, \(\eta\) 为 \(Y\) 的极大点,
则 \(\eta\) 的各闭邻域的交等于 \(\overline{\{\eta\}}\).
事实上, 若 \(y\in Y-\overline{\{\eta\}}\),
则 \(y\) 属于某个不包含 \(\eta\) 的仿射开集 \(V\).
由于 \(Y\) 拟分离, \(V\) 相对拟紧,
因而 \(\overline V\) 为 \(\overline{\{\xi\}}\) 之并,
其中 \(\xi\) 遍历 \(Y\) 中属于 \(V\) 的极大点.
故 \(\eta\notin\overline V\), 因而 \(Y-\overline V\)
是 \(\eta\) 的不包含 \(y\) 的闭邻域.

\medskip
\phantomsection\label{prop:sga3-VIA-2.6.4}%
\noindent\textbf{命题 2.6.4.}
\textit{设 \(X\) 为满足 \((\star)\) 的 \(G\)-概形, \(U\) 为 \(X\) 的开子集.}
\begin{enumerate}
\item[\SGAInlineRefTarget{prop:sga3-VIA-2.6.4:item:i}(i)] \textit{\(C^0U\) 为 \(X\) 的开子集, 等于 \(X\) 中泛点属于 \(U\) 的不可约分支之并.}
\item[(i\(^{\prime}\))] \textit{因此, 若 \(X\) 不可约, 则它拟紧.}
\item[\SGAInlineRefTarget{prop:sga3-VIA-2.6.4:item:ii}(ii)] \textit{此外, 若 \(U\) 在 \(X\) 中相对拟紧,
则 \(C^0U=\overline U\), 且此子集在 \(X\) 中既开又闭.}
\end{enumerate}

\noindent\textit{证明.}
\SGARefLink{prop:sga3-VIA-2.6.4:item:i}{(i)} 首先, \(C^0U\) 为开集, 因为它是如下开子集的投影之并,
其中 \(g\in C^0\):
\[
  g\cdot U_{\kappa(g)}\subset X_{\kappa(g)},
\]
而每个投影 \(X_{\kappa(g)}\to X\) 都是开映射.

为证明\SGARefLink{prop:sga3-VIA-2.6.4:item:i}{(i)} 的第二个断言,
可将 \(k\) 换为 \(k'\) (因为 \(\operatorname{Spec}k'\to\operatorname{Spec}k\)
为普遍同胚), 从而假设 \(k\) 完美.
进而可假设 \(G\) 和 \(X\) 约化, 从而几何约化.

设 \(\eta\) 为 \(X\) 中包含于 \(U\) 的极大点,
令 \(Z=\overline{\{\eta\}}\). 考虑
\[
  \mu':C^0\times Z\longrightarrow X.
\]
由于 \(C^0\) 几何不可约, \(C^0\times Z\) 不可约;
用 \(\gamma\) 表示其泛点. 态射 \(\mu'\) 将点 \(\varepsilon(\eta)\)
映到 \(\eta\), 其中 \(\varepsilon\) 表示 \(C^0\) 的单位截面.
故 \(\gamma\) 被映到 \(\eta\) 的一个泛化点, 从而被映到 \(\eta\).
因此 \(\mu'\) 将 \(C^0\times Z\) 的底空间映入 \(Z\),
且由于 \(C^0\times Z\) 约化, \(\mu'\) 经由 \(Z\) 分解.

现在设 \(z\in Z\), 并令 \(K=\kappa(z)\). 态射
\[
  \mu_z:G_K\longrightarrow X_K,\qquad h\longmapsto h\cdot z,
\]
为满射. 设 \(\alpha\) 为 \(X_K\) 的极大点.
由于 \(X_K\) 约化, 局部环 \(\mathcal O_{X_K,\alpha}\) 为域,
故 \(\mu_z\) 在 \(G_K\) 中位于 \(\alpha\) 上方的每个点平坦.
由\hyperref[lem:sga3-VIA-2.5.3]{引理~2.5.3} 可知 \(\mu_z\) 平坦.

另一方面, \(\mu_z\) 将 \(C^0_K\) 的泛点 \(\omega\) 映到某点 \(t\in Z_K\).
设 \(\beta\) 为 \(Z_K\) 的极大点, 使 \(t\in\overline{\{\beta\}}\).
由于 \(\mu_z\) 平坦, 存在 \(\omega\) 的泛化点 \(\xi\),
使 \(\mu_z(\xi)=\beta\). 但 \(\omega\) 是 \(G_K\) 的极大点,
故必有 \(\xi=\omega\). 因此 \(\mu_z(\omega)=\beta\),
且由于 \(Z_K\to Z\) 平坦, \(\beta\) 位于 \(\eta\) 上方.

令 \(L=\kappa(\omega)\), 并设 \(\omega_L\) 和 \(z_L\)
为由 \(\omega\) 及 \(z\) 得到的 \(L\)-点. 则
\[
  \omega_L\cdot z_L=\beta'
\]
为 \(Z_L\) 中位于 \(\beta\) 上方的一点, 因而
\[
  z_L=\omega_L^{-1}\cdot\beta'\in C^0_L\cdot U_L.
\]
因此 \(z\in C^0\cdot U\), 这就证明了\SGARefLink{prop:sga3-VIA-2.6.4:item:i}{(i)}.

由\hyperref[cor:sga3-VIA-2.6.2]{2.6.2}, \(C^0\) 拟紧,
从而得到断言 (i\(^{\prime}\)): 若 \(X\) 不可约,
\(U\) 为非空仿射开子集, 则 \(X=C^0U\),
故 \(X\) 为态射 \(C^0\times U\to X\) 的像, 因而拟紧.
最后, 若 \(U\) 在 \(X\) 中相对拟紧,
则由\hyperref[rem:sga3-VIA-2.6.3]{2.6.3}, \(\overline U\)
是 \(\overline{\{\eta\}}\) 之并, 其中 \(\eta\) 遍历 \(X\) 中包含于 \(U\) 的极大点.
故 \(\overline U=C^0U\). 这就证明了\SGARefLink{prop:sga3-VIA-2.6.4:item:ii}{(ii)}.
% Independent prose review: source and translation compared paragraph by paragraph through this component; statements, hypotheses, proofs, notes and diagram captions covered; bibliography reserved for its assigned reviewer; source math and reference keys unchanged.
